\subsection{Fixed finite local matching and a later metabelian obstruction}

\begin{statement}{A.202}{Proved}
{Ordinary ghosts with a prescribed nonzero dyadic separation depth}
Let \(N\triangleleft\Pi\) be open and normal in \(E_U\), put
\(Q=\Pi/N\), and use the actual quotient
\[
 E_{Q,A}=E_U/(N\cap\overline{[\Pi,\Pi]}),
 \qquad A=\Pi^{\rm ab},
\tag{A.202.1}
\]
with its point-three section in the \(Q\)-factor.
For every sufficiently large integer \(E\), there is a continuous
section \(\tau\) of (A.202.1) with an ordinary infinite-support
completed Kummer pair \((A_\tau,B_\tau)\) such that
\[
 \begin{aligned}
 (\kappa_8(A_\tau),\kappa_8(B_\tau))&=([3],[-2]),\\
 v_2(x_2-3)&=E,
 \end{aligned}
\tag{A.202.2}
\]
where \((x_2,1-x_2)\) is its ordinary dyadic localization.
Its \(Q\)-projection is the point-three section, and at every
finite place \(\tau\) is geometrically conjugate to the image
of an ordinary local point.  At the real place it is the image
of the point three.

The same completed pair has an actual global lift to the full
profinite class-two arithmetic quotient by A.38.  It has no
everywhere locally point-or-cusp-matching lift to the full
metabelian arithmetic quotient \(E_U/\Pi''\), and consequently
none to a full Selmer section of \(U\).
These are separate lifting assertions: no joint locally matching
lift retaining both \(Q\) and the class-two image is asserted.

The obstruction has an actual finite target.  For the cyclic-eight
cover \(W\to U\) of A.192, let \(E_n\) be the arithmetic
quotient in (A.193.13).  There is no everywhere locally
point-or-cusp-matching section of
\[
 E_{E-2},\qquad
 |\ker(E_{E-2}\to G_\Q)|=2^{4E-5},\qquad
 \operatorname{cl}\ker(E_{E-2}\to G_\Q)=4E-8,
\tag{A.202.3}
\]
with combined Kummer class \([-24]\) and ordinary dyadic
centre \(x_2\).  This finite quotient is chosen after the
dyadic depth has been prescribed; the initial \(Q\) is arbitrary
and fixed before the construction.
\end{statement}

\begin{proof}
\emph{The literal change to the prime recursion.}
Keep the finite cover, splitting field, constants \(C,N_0\),
exponents \(n_j=N_0(j+4)!\), and the distinguished components
\(W_j\) of A.72.  In particular \(8\mid n_j\).
Write the evolving coefficients as
\[
 a_0=3,\quad b_0=-2,\qquad
 a_{j+1}=a_jq_j^{n_j},\quad b_{j+1}=b_jr_j^{n_j}.
\tag{A.202.4}
\]
At two, the controlled local point in (A.72.11) has first
Kummer coordinate \(\xi_2=1\) at every stage.  This equality
is retained by the actual transition (A.72.18).

At the initial stage choose \(L\) at least as large as the
dyadic precision required in (A.72.13)--(A.72.15), and also
\(L\geq\max(2,v_2(n_1))\).  Replace only the first dyadic
prime congruence by
\[
                    q_0\equiv1+2^L\pmod {2^{L+1}}.
\tag{A.202.5}
\]
This still implies the required \(q_0\equiv\xi_2=1\pmod{2^L}\).
All other congruences and splitting conditions of (A.72.16)
are kept.  Their combined CRT class is one modulo \(n_0\).
The proved field intersection (A.72.4), now with the dyadic
part of the auxiliary modulus enlarged to \(2^{L+1}\),
therefore gives exactly the required Chebotarev element.
It supplies a fresh prime \(q_0\) satisfying (A.202.5).
The second prime \(r_0\) is chosen as in (A.72.17), including
its old dyadic precision; that argument is unchanged.

For \(q\equiv1\pmod4\), repeated factorization of
\(q^{2h}-1=(q^h-1)(q^h+1)\), followed by the odd-exponent
geometric sum, proves
\[
             v_2(q^n-1)=v_2(q-1)+v_2(n).
\tag{A.202.6}
\]
Thus the first updated coefficient has
\[
       v_2(a_1-3)=L+v_2(n_0)=:E.
\tag{A.202.7}
\]
Every sufficiently large prescribed \(E\) is obtained by
taking \(L=E-v_2(n_0)\); only the fixed initial lower bounds
on \(L\) have to be met.

At each later stage increase the dyadic precision \(L_j\)
as necessary so that, in addition to all the original local
continuation and Hensel inequalities,
\[
              L_j\geq2,\qquad L_j+v_2(n_j)\geq E+1.
\tag{A.202.8}
\]
Keep \(q_j\equiv1\pmod {2^{L_j}}\).  Equation (A.202.6)
implies \(q_j^{n_j}\in1+2^{E+1}\Z_2\).
Larger precisions are allowed at every step of A.72, and its
same intersection calculation proves all prime choices remain
available.  No local point is moved to a different component:
(A.72.14)--(A.72.18) use only the retained congruences to the
old point and become stronger when the precision increases.
Consequently the original all-place unit-point invariant,
the exact cover maps and the ordinary limits all remain valid.

The later product converges in the closed subgroup
\(1+2^{E+1}\Z_2\).  The ordinary dyadic limit therefore obeys
\[
 x_2=3q_0^{n_0}\prod_{j\geq1}q_j^{n_j},\qquad
 x_2\equiv a_1\pmod {2^{E+1}},\qquad
                       v_2(x_2-3)=E.
\tag{A.202.9}
\]
The first congruence in (A.202.2) follows directly from
\(8\mid n_j\) and the cofinal convergence of the completed
products.  More explicitly, all tails reduce identically to
one at exponent eight.  The two disjoint fresh-prime supports
remain infinite.  A.72's actual fibre-product section,
finite-cover local fixed sheets and compact conjugator
intersection now give all the asserted joint local matching
in \(E_{Q,A}\), including the real point-three image.

\emph{The complete dyadic cover calculation.}
The point-three frame makes
\(t(1-t)^3/(-24)\) an eighth power at \(t=x_2\).
Thus its actual torsor \(W\) has a dyadic sheet.
Every such sheet is an ordinary point of the affine genus-two
curve \(C\), with coordinates which we denote \((x,y)\)
to distinguish its coordinate from \(x_2\in U(\Q_2)\).
The exact map of A.192 is
\[
 t=\frac{x^4}{x^4-54},\qquad
 x^4=\frac{54t}{t-1}.
\tag{A.202.10}
\]
Since \(E\) is large, \(t=x_2\equiv3\pmod8\), so
\(v_2(t)=0\) and \(v_2(t-1)=1\).  The second equality
in (A.202.10) gives \(v_2(x)=0\).
For an odd \(x\), the equation \(y^2=x^5-54x\) reduces
modulo eight to \(y^2\equiv3x\).  Its right side is odd,
so the square condition gives \(x\equiv3\pmod8\).

The exact factorization is
\[
 t-3=-\frac{2(x-3)(x+3)(x^2+9)}{x^4-54}.
\tag{A.202.11}
\]
For \(x\equiv3\pmod8\), the denominator is a unit,
\(v_2(x+3)=1\), and \(v_2(x^2+9)=1\), the last assertion
following from \(x^2\equiv9\pmod {16}\).
Since \(t\ne3\), we have \(x\ne3\).  Hence on every sheet
\[
              v_2(t-3)=v_2(x-3)+3,
                  \qquad v_2(x-3)=E-3.
\tag{A.202.12}
\]
This also follows from the literal deck action:
\(\mu_8(\Q_2)=\{1,-1\}\), and the inverse formulas
\(x=-6z^2/(1-t)\), \(y=-18z/(1-t)\) keep \(x\) fixed
and negate \(y\).  Thus no choice of the fixed sheet avoids
the displayed valuation.

An everywhere locally point-or-cusp-matching section of \(E_n\)
in the combined class \([-24]\) has a fixed sheet on \(W\).
Compactifying its cover subgroup gives the actual global
class in \(\operatorname{Sel}^{2^n}(J/\Q)\) from
(A.193.13), with the finite local Kummer images of the selected
proper centres.  The last row of (A.193.3) is the exact condition
\(n\leq v_2(x-3)\).  Setting \(n=E-2\) contradicts
(A.202.12).  The geometric order and class in (A.202.3)
are the substitutions into (A.193.4).  This proves the finite
obstruction with its actual arithmetic map and every sheet.

Finally suppose a locally matching full metabelian section had
the completed pair constructed above.  At two its point or cusp
image has the same two full local Kummer classes as
\((x_2,1-x_2)\).  The actual local Kummer injection fixes an
ordinary centre to be \(x_2\).  It also excludes all three
cusp packets: at zero the second class is one, contrary to
\(v_2(1-x_2)=1\); at one the first class is one, contrary
to \(x_2\equiv3\pmod8\); and at infinity the two valuation
coordinates are equal, contrary to \((0,1)\).
Its projection through \(\Pi''\subset\ker(\Pi\to Q_n)\)
would now give the forbidden section of \(E_{E-2}\).
This proves the asserted nonexistence, also for full Selmer
sections through the actual projection to the metabelian quotient.

On the other hand the constructed pair is in the ordinary
completed unit-equation locus \(\mathcal G\), so A.38 gives
its global continuous class-two lift through
\(E_U/\Gamma_3\Pi\to E_U/\overline{[\Pi,\Pi]}\).
That theorem specifies a global lift of this abelian frame;
the local matching proved here is in (A.202.1).
No simultaneous local conjugator or common lift of those two
quotients is used or asserted.

The retained certificate
\path{code/point3_fixed_quotient_separation_certificate.py}
checks the literal depth identities, the allowed precision
increases, the cover function and deck identities, and the
finite detector index/order/class.  The infinite prime recursion
and all-place matching use the proved exact modifications of
A.72 above; no large first prime search is claimed as executed.
\end{proof}
